\documentclass[12pt,a4paper]{article} %deklarace typu dokumentu
\usepackage[utf8]{inputenc} %kódování
\usepackage[czech]{babel} %jazyk dokumentu
\usepackage{fixltx2e}
\usepackage[unicode]{hyperref}
\usepackage{siunitx}
\usepackage{graphicx}
%rovnice k levému okraji
\usepackage{amsmath}
\graphicspath{ {./images/} }

\title{Semestrální projekt - Elektronika pro informační technologie}
\author{Pavel Šesták [xsesta07]}
\date{3.10.2019}                                         
                                              
\begin{document}
    \begin{titlepage}
		\begin{center}
			\textsc{\LARGE Vysoké Učení Technické v Brně} \\[0.5cm]
			{\LARGE Fakulta informačních technologií}

			\begin{figure}
				\center\includegraphics[scale=0.3]{images/logo.png}
			\end{figure}

			\vspace{3cm}

			\textsc{\LARGE ÚVOD DO SOFTWAROVÉHO INŽENÝRSTVÍ} \\[0.5cm]
			\textsc{\LARGE 2019/2020} \\[3.5cm]

			\textbf{\LARGE Projekt - Model informačního systému} \\[2cm]
			\textbf{Zadání č. 4 – Prodejna se sýry}
		\end{center}
		\vfill
		\begin{flushleft} 
			\large
			Pavel Šesták (xsesta07) 
			\hfill
			Brno, \today
		\end{flushleft}
	\end{titlepage}
    
    \section*{Zadání -  4. Prodejna se sýry}
   
Navrhněte IS prodejny se sýry, který by měl umožňovat evidovat zboží v prodejně a v jejím skladu. Dále by měl umožnit zaměstnanci objednávat sýry v prodejně a na skladu od různých dodavatelů. Jedna objednávka může obsahovat i více druhů zboží. Sýry se kategorizují podle více kritérií: země původu, živočicha (kravský, ovčí, kozí, jačí, ...), typu (ementálový, plísňový, uzený, ...). U sýru je nutno evidovat procento tuku. Sýry jsou do prodejny dodávány v bochnících, kde pro každý bochník je třeba evidovat jeho počáteční hmotnost, aktuální hmotnost na prodejně, datum dodání, trvanlivost. Zaměstnanci si mohou vyhledávat sýry i podle země původu, u které jsou k dispozici relevantní informace o sýrech z dané země, či podle jeho druhu nebo živočicha. Zaměstnanci prodejny mohou objednávat různé druhy sýrů u různých dodavatelů (ne všichni dodavatelé dodávají všechny druhy sýra), přičemž objednávka může obsahovat více druhů sýrů a u každého druhu je nutno specifikovat množství (hmotnost). Pro dodané sýry je nutno mít možnost dohledat, ke které objednávce patří a od kterého dodavatele jsou.
    \clearpage
    \section*{Use case} % nečíslovaný nadpis, není uveden v obsahu
        \center\includegraphics[scale=0.5]{images/IUS_semestralka_UseCaseDia.png}
	\clearpage
	\section*{ERD Diagram} % nečíslovaný nadpis, není uveden v obsahu
        
        %\center\includegraphics[scale=0.45]{images/ERD_diagram.png}

            \includegraphics[width=1.3\textwidth, angle =90 ]{images/ERD_diagram.png}

	\clearpage


	


\end{document}
